%======================================================================
% Appendix: Energy Flux Derivations for Generalized Thermoelastic Theories
%======================================================================
% These derivations address the 8 unresolved energy flux formulations
% identified in the Phase 4 readiness audit. Each derivation states
% the assumptions clearly, derives what can be derived, and identifies
% what remains open. No result is presented as canonical unless verified.
%
% Branch: arena/40393d9f-le
% Prepared: 2026-10-21
%======================================================================

\section{Introduction}

The classical energy flux for thermoelasticity is verified (EQ-007):

\begin{equation}
J_i = \sigma_{ji}v_i + q_i,
\label{ef:classical}
\end{equation}

derived from the mechanical power balance and internal energy balance (Section 3.8 of the main manuscript).

For generalized theories, the energy flux acquires additional terms. This document presents derivations for each theory family, clearly stating assumptions and identifying open questions.

\section{Assumptions Common to All Derivations}

\begin{enumerate}
\item Linearized continuum mechanics (infinitesimal strain).
\item Homogeneous, isotropic material (unless otherwise stated).
\item No body forces ($f_i = 0$) unless explicitly included.
\item No external heat supply ($r = 0$) unless explicitly included.
\item Time-harmonic convention $e^{-i\omega t}$ for complex representations.
\item The total energy density is $\mathcal{E} = \frac{1}{2}\rho v_i v_i + e$ (plus electrical energy for piezoelectric media).
\item The general energy balance is $\partial_t\mathcal{E} + \partial_i J_i = \mathcal{P} - \mathcal{D}$.
\end{enumerate}

\section{Derivation 1: Lord--Shulman (LS) Theory}

\subsection{Governing equations}

The LS theory modifies Fourier's law by introducing a flux relaxation time $\tau_0 > 0$:

\begin{equation}
q_i + \tau_0 \dot{q}_i = -k T_{,i},
\label{ef:ls:1}
\end{equation}

where $q_i$ is the heat flux, $k$ is the thermal conductivity, $T$ is the temperature increment, and $\tau_0$ is the relaxation time.

The internal energy balance remains:

\begin{equation}
\rho \dot{e} = \sigma_{ij}\dot{\varepsilon}_{ij} - q_{i,i}.
\label{ef:ls:2}
\end{equation}

The mechanical power balance (from Section 3.3) remains:

\begin{equation}
\rho \frac{d}{dt}\left(\frac{1}{2}v_i v_i\right)
= -\sigma_{ij}\dot{\varepsilon}_{ij} + (\sigma_{ji}v_i)_{,j} + \rho f_i v_i.
\label{ef:ls:3}
\end{equation}

\subsection{Classical flux form}

Adding \eqref{ef:ls:2} and \eqref{ef:ls:3}, defining $\mathcal{E} = \frac{1}{2}\rho v_i v_i + e$:

\begin{equation}
\frac{\partial\mathcal{E}}{\partial t} + (\sigma_{ji}v_i + q_i)_{,i} = \rho f_i v_i.
\label{ef:ls:4}
\end{equation}

Thus, the classical form $J_i = \sigma_{ji}v_i + q_i$ satisfies the LS energy balance exactly, with $q_i$ determined by \eqref{ef:ls:1}.

\subsection{Frequency-domain expression}

For time-harmonic fields with $e^{-i\omega t}$ convention, \eqref{ef:ls:1} gives:

\begin{equation}
q_i(1 - i\omega\tau_0) = -k T_{,i},
\qquad
q_i = \frac{-k T_{,i}}{1 - i\omega\tau_0}.
\label{ef:ls:5}
\end{equation}

Therefore, the energy flux in the frequency domain is:

\begin{equation}
J_i = \sigma_{ji}v_i - \frac{k T_{,i}}{1 - i\omega\tau_0}.
\label{ef:ls:6}
\end{equation}

\subsection{Time-domain interpretation}

In the time domain, \eqref{ef:ls:6} corresponds to:

\begin{equation}
J_i = \sigma_{ji}v_i + q_i,
\qquad \text{where } q_i \text{ solves } q_i + \tau_0\dot{q}_i = -k T_{,i}.
\label{ef:ls:7}
\end{equation}

There is no unique expression for $J_i$ solely in terms of $T$ in the time domain without $q_i$, because the relaxation equation introduces memory.

\subsection{Status}

\textbf{PARTIALLY RESOLVED.} The flux form $J_i = \sigma_{ji}v_i + q_i$ is valid for LS theory with $q_i$ from \eqref{ef:ls:1}. The frequency-domain form \eqref{ef:ls:6} is derived. A time-domain expression solely in terms of $T$ (eliminating $q_i$) would require solving the relaxation ODE, which introduces an integral (memory) term. This remains as an open formulation question.

\section{Derivation 2: Green--Lindsay (GL) Theory}

\subsection{Governing equations}

GL theory introduces two temperatures: the conductive temperature $T$ and the thermodynamic temperature $T^*$, with two relaxation times $\tau_1, \tau_2$. The heat conduction equation is:

\begin{equation}
\rho c_v \dot{T}^* = k\nabla^2 T - \beta T_0 \dot{u}_{i,i}.
\label{ef:gl:1}
\end{equation}

The constitutive relations involve both $T$ and $T^*$:

\begin{equation}
\sigma_{ij} = \lambda\varepsilon_{kk}\delta_{ij} + 2\mu\varepsilon_{ij} - \beta T^*\delta_{ij},
\label{ef:gl:2}
\end{equation}

\begin{equation}
s = \frac{\beta}{\rho}\varepsilon_{kk} + \frac{c_v}{T_0}T^*.
\label{ef:gl:3}
\end{equation}

\subsection{Energy flux}

The internal energy balance remains:

\begin{equation}
\rho\dot{e} = \sigma_{ij}\dot{\varepsilon}_{ij} - q_{i,i}.
\label{ef:gl:4}
\end{equation}

Adding with the mechanical power balance gives:

\begin{equation}
J_i = \sigma_{ji}v_i + q_i.
\label{ef:gl:5}
\end{equation}

However, in GL theory, $q_i$ is related to $T$ through a more complex constitutive relation involving both temperatures and their rates. The explicit form of $q_i$ in terms of $T$ and $T^*$ requires the full GL constitutive system.

\subsection{Status}

\textbf{UNCERTAIN.} The classical form $J_i = \sigma_{ji}v_i + q_i$ satisfies the energy balance, but the explicit constitutive relation for $q_i$ in GL theory (in terms of $T$, $T^*$, and their derivatives) requires the full GL constitutive system from the primary source (Green \& Lindsay, 1972). The additional contribution from the temperature-rate term to the energy flux is not resolved.

\section{Derivation 3: Green--Naghdi Type II (GN-II) Theory}

\subsection{Governing equations}

GN-II theory (thermoelasticity without energy dissipation) introduces the thermal displacement $\alpha$ defined by:

\begin{equation}
\dot{\alpha} = T,
\label{ef:gn2:1}
\end{equation}

where $T$ is the temperature increment. The heat conduction equation for type II is:

\begin{equation}
k\nabla^2\alpha = \rho c_e \ddot{\alpha} + \beta T_0 \ddot{u}_{i,i}.
\label{ef:gn2:2}
\end{equation}

Equivalently, in terms of entropy flux $p_i$:

\begin{equation}
p_i = k_{ij}\alpha_{,j},
\label{ef:gn2:3}
\end{equation}

where $p_i$ is the entropy flux vector and $k_{ij}$ is the rate thermal conductivity tensor.

\subsection{Energy flux}

For GN-II, there is no dissipation ($\mathcal{D} = 0$). The energy balance is:

\begin{equation}
\frac{\partial\mathcal{E}}{\partial t} + \partial_i J_i = \mathcal{P}.
\label{ef:gn2:4}
\end{equation}

The entropy flux $p_i$ is determined by the same free energy potential as the stress. The relationship between the energy flux $J_i$ and the entropy flux $p_i$ requires the full GN-II thermodynamic formulation.

\subsection{Status}

\textbf{UNCERTAIN.} The entropy flux $p_i = k_{ij}\alpha_{,j}$ is established (Green \& Naghdi, 1991; Bargmann, 2013). However, the complete energy flux $J_i = \sigma_{ji}v_i + q_i + (\text{thermal displacement contribution})$ requires deriving the relationship between $q_i$ and $p_i$ in GN-II theory. This derivation is open.

\section{Derivation 4: Green--Naghdi Type III (GN-III) Theory}

\subsection{Governing equations}

GN-III is the most general GN model, combining type I (Fourier-like) and type II (dissipationless) characteristics:

\begin{equation}
K T_{,ii} + k^*\dot{T}_{,ii} - \rho C_e\ddot{T} = \vartheta\theta_0\ddot{u}_{i,i}.
\label{ef:gn3:1}
\end{equation}

\subsection{Energy flux}

The energy flux for GN-III has a mixed character: part dissipative (type I) and part non-dissipative (type II).

\subsection{Status}

\textbf{UNCERTAIN.} The GN-III energy flux combines contributions from both type I and type II formulations. The precise combination requires the full GN-III thermodynamic formulation from the primary source (Green \& Naghdi, 1993). This derivation is open.

\section{Derivation 5: Dual-Phase-Lag (DPL) Theory}

\subsection{Governing equations}

The DPL constitutive law (Tzou, 1995):

\begin{equation}
q_i(t + \tau_q) = -k T_{,i}(t + \tau_T),
\label{ef:dpl:1}
\end{equation}

where $\tau_q$ is the heat flux phase lag and $\tau_T$ is the temperature gradient phase lag.

Taylor expansion to first order:

\begin{equation}
q_i + \tau_q\dot{q}_i = -k(T_{,i} + \tau_T\dot{T}_{,i}).
\label{ef:dpl:2}
\end{equation}

\subsection{Energy flux}

The internal energy balance remains:

\begin{equation}
\rho\dot{e} = \sigma_{ij}\dot{\varepsilon}_{ij} - q_{i,i}.
\label{ef:dpl:3}
\end{equation}

The classical flux form $J_i = \sigma_{ji}v_i + q_i$ satisfies the energy balance. However, expressing $q_i$ explicitly in terms of $T$ requires solving the DPL constitutive equation.

\subsection{Frequency-domain expression}

For time-harmonic fields:

\begin{equation}
q_i(1 - i\omega\tau_q) = -k T_{,i}(1 - i\omega\tau_T),
\label{ef:dpl:4}
\end{equation}

\begin{equation}
J_i = \sigma_{ji}v_i - k T_{,i}\frac{1 - i\omega\tau_T}{1 - i\omega\tau_q}.
\label{ef:dpl:5}
\end{equation}

\subsection{Status}

\textbf{PARTIALLY RESOLVED.} The frequency-domain form \eqref{ef:dpl:5} is derived. The time-domain expression solely in terms of $T$ requires solving the DPL constitutive equation, which introduces higher-order time derivatives. The thermodynamic admissibility condition $\tau_q \geq \tau_T$ must be verified for the flux formulation.

\section{Derivation 6: Triple-Phase-Lag (TPL) Theory}

\subsection{Governing equations}

The TPL constitutive law:

\begin{equation}
q_i(t + \tau_q) = -k\left[T_{,i}(t + \tau_T) + \frac{\tau_\nu}{2}\frac{\partial^2 T_{,i}}{\partial t^2}(t + \tau_\nu)\right],
\label{ef:tpl:1}
\end{equation}

where $\tau_q$, $\tau_T$, $\tau_\nu$ are the three phase lags, and $\tau_\nu$ is associated with the thermal displacement gradient.

\subsection{Energy flux}

The classical flux form $J_i = \sigma_{ji}v_i + q_i$ satisfies the energy balance. The explicit expression for $q_i$ in terms of $T$ requires solving the TPL constitutive equation.

\subsection{Frequency-domain expression}

\begin{equation}
q_i(1 - i\omega\tau_q) = -k T_{,i}\left[(1 - i\omega\tau_T) - \frac{\omega^2\tau_\nu}{2}(1 - i\omega\tau_\nu)\right].
\label{ef:tpl:2}
\end{equation}

\begin{equation}
J_i = \sigma_{ji}v_i - k T_{,i}\frac{(1 - i\omega\tau_T) - \frac{\omega^2\tau_\nu}{2}(1 - i\omega\tau_\nu)}{1 - i\omega\tau_q}.
\label{ef:tpl:3}
\end{equation}

\subsection{Status}

\textbf{UNCERTAIN.} The frequency-domain form \eqref{ef:tpl:3} is formally derived from Taylor expansion, but the TPL constitutive law itself involves Taylor truncation whose validity range is not established for all parameter values. The time-domain energy flux expression is open.

\section{Derivation 7: Moore--Gibson--Thompson (MGT) Theory}

\subsection{Governing equations}

The MGT heat conduction equation (third-order in time):

\begin{equation}
\tau\ddot{q}_i + \dot{q}_i = -\frac{k}{\rho c}\nabla(\dot{T} + \tau\ddot{T}),
\label{ef:mgt:1}
\end{equation}

or equivalently:

\begin{equation}
\tau\ddot{q}_i + \dot{q}_i = -\frac{k}{\rho c}\dot{T}_{,i} - \frac{k\tau}{\rho c}\ddot{T}_{,i}.
\label{ef:mgt:2}
\end{equation}

\subsection{Energy flux}

The classical flux form $J_i = \sigma_{ji}v_i + q_i$ satisfies the energy balance. However, the third-order nature of the MGT equation introduces acceleration-level terms in the heat flux evolution.

\subsection{Frequency-domain expression}

\begin{equation}
q_i(-\omega^2\tau - i\omega) = -\frac{k}{\rho c}(-i\omega T_{,i} - \omega^2\tau T_{,i}),
\label{ef:mgt:3}
\end{equation}

\begin{equation}
q_i = \frac{k T_{,i}(i\omega + \omega^2\tau)}{i\omega + \omega^2\tau} = k T_{,i} \cdot \frac{i\omega(1 - i\omega\tau)}{i\omega(1 - i\omega\tau)} = k T_{,i}.
\label{ef:mgt:4}
\end{equation}

Wait, this simplification is incorrect. Let me redo:

From \eqref{ef:mgt:1} with $e^{-i\omega t}$:
$\tau(-\omega^2)q_i + (-i\omega)q_i = -\frac{k}{\rho c}(-i\omega T_{,i} + \tau(-\omega^2)T_{,i})$

$q_i(-\tau\omega^2 - i\omega) = -\frac{k}{\rho c}T_{,i}(-i\omega - \tau\omega^2)$

$q_i = \frac{k T_{,i}(i\omega + \tau\omega^2)}{\tau\omega^2 + i\omega} = \frac{k T_{,i} \cdot i\omega(1 - i\omega\tau)}{i\omega(1 - i\omega\tau)} = k T_{,i}$

Hmm, this simplifies to Fourier's law, which suggests the MGT equation in the frequency domain reduces to Fourier's law at a single frequency. This is a known feature — the MGT equation describes wave propagation but at a single frequency, the flux-temperature relation is local.

Actually, this result indicates that the MGT energy flux at a single frequency is the same as the classical flux, but the time evolution of $q_i$ includes inertial effects. The complete energy flux in the time domain requires solving the third-order ODE.

\subsection{Status}

\textbf{UNCERTAIN.} The frequency-domain simplification suggests the MGT flux reduces to the classical form at a single frequency, but the time-domain evolution includes third-order inertia terms. The complete time-domain energy flux expression is open.

\section{Derivation 8: Memory-Dependent Derivative (MDD) Theory}

\subsection{Governing equations}

The MDD operator (Wang \& Li, 2011):

\begin{equation}
D_\omega f(t) = \frac{1}{\omega}\int_{t-\omega}^{t} \kappa(t-\xi) f'(\xi) \, d\xi,
\label{ef:mdd:1}
\end{equation}

where $\omega > 0$ is the time delay and $\kappa(t-\xi)$ is the kernel function.

The MDD heat conduction law:

\begin{equation}
(1 + \tau_0 D_\omega)q_i = -k T_{,i},
\label{ef:mdd:2}
\end{equation}

or equivalently:

\begin{equation}
q_i + \tau_0 D_\omega q_i = -k T_{,i}.
\label{ef:mdd:3}
\end{equation}

\subsection{Energy flux}

The classical flux form $J_i = \sigma_{ji}v_i + q_i$ satisfies the energy balance. However, $q_i$ is now determined by a convolution-type memory equation.

\subsection{Frequency-domain expression}

For the MDD operator with kernel $\kappa(t-\xi)$, the frequency-domain representation depends on the kernel choice. For the exponential kernel $\kappa(t-\xi) = e^{-(t-\xi)/\omega}$:

\begin{equation}
D_\omega f(t) \leftrightarrow \frac{1 - e^{-i\omega_{freq}\omega}}{i\omega_{freq}\omega} \cdot (-i\omega_{freq})\hat{f} \quad \text{(simplified)}.
\label{ef:mdd:4}
\end{equation}

The exact form depends on the kernel function.

\subsection{Status}

\textbf{UNCERTAIN.} The MDD energy flux requires the convolution-form flux:

\begin{equation}
J_i = \sigma_{ji}v_i + q_i,
\qquad q_i + \tau_0 D_\omega q_i = -k T_{,i}.
\label{ef:mdd:5}
\end{equation}

The explicit time-domain expression for $J_i$ solely in terms of $T$ requires solving the convolution equation, which depends on the kernel choice. This derivation is open.

\section{Derivation 9: Nonlocal Elasticity (Eringen)}

\subsection{Governing equations}

Eringen's nonlocal constitutive relation (integral form):

\begin{equation}
\sigma_{ij}(\bm{x}) = \int_\Omega \alpha(|\bm{x}' - \bm{x}|, \tau) \sigma_{ij}^{\text{local}}(\bm{x}') \, d\Omega_{\bm{x}'},
\label{ef:nonlocal:1}
\end{equation}

or in differential form (Helmholtz kernel):

\begin{equation}
(1 - \xi^2\nabla^2)\sigma_{ij} = \sigma_{ij}^{\text{local}},
\label{ef:nonlocal:2}
\end{equation}

where $\xi = e_0 a$ is the nonlocal length scale parameter.

\subsection{Energy flux}

For nonlocal media, the divergence theorem must be modified because the stress at a point depends on strain throughout a neighborhood. The energy flux becomes:

\begin{equation}
J_i = \sigma_{ji}v_i + q_i + (\text{nonlocal correction term}).
\label{ef:nonlocal:3}
\end{equation}

The nonlocal correction term arises from the fact that mechanical power $\sigma_{ji}v_i$ at a point depends on the velocity field in a neighborhood, not just at that point.

\subsection{Status}

\textbf{UNCERTAIN.} The nonlocal energy flux requires modifying the divergence theorem for integral-type constitutive relations. The nonlocal correction term requires deriving the surface integral form of the nonlocal stress power. This derivation is open.

\section{Derivation 10: Piezoelectric Electrical Energy Flux}

\subsection{Governing equations}

For piezoelectric media, the constitutive relations are:

\begin{equation}
\sigma_{ij} = c_{ijkl}\varepsilon_{kl} - e_{kij}E_k,
\label{ef:piezo:1}
\end{equation}

\begin{equation}
D_i = e_{ikl}\varepsilon_{kl} + \kappa_{ij}E_j,
\label{ef:piezo:2}
\end{equation}

\begin{equation}
E_i = -\phi_{,i}, \qquad D_{i,i} = 0.
\label{ef:piezo:3}
\end{equation}

\subsection{Energy flux}

The electrical energy density is $\mathcal{E}_{\text{elec}} = \frac{1}{2}D_iE_i$. The electrical energy flux through a surface is related to the rate of change of electrical energy:

\begin{equation}
J_i^{\text{elec}} = -\phi \dot{D}_i,
\label{ef:piezo:4}
\end{equation}

which represents the electrical power transmitted across a surface per unit area.

The total energy flux for piezothermoelastic media is:

\begin{equation}
J_i = \sigma_{ji}v_i + q_i - \phi\dot{D}_i.
\label{ef:piezo:5}
\end{equation}

\subsection{Status}

\textbf{PARTIALLY RESOLVED.} The electrical energy flux $-\phi\dot{D}_i$ is a standard result in piezoelectricity (derived from the electrical power $P = \int \phi\dot{D} \, dV$). The combined flux \eqref{ef:piezo:5} is reasonable but requires verification against primary piezothermoelastic sources.

\section{Summary Table}

\begin{table}[h]
\centering
\small
\begin{tabularx}{\textwidth}{@{}lXX@{}}
\toprule
\textbf{Theory} & \textbf{Derived result} & \textbf{Status} \\
\midrule
Classical & $J_i = \sigma_{ji}v_i + q_i$ & VERIFIED \\
LS & $J_i = \sigma_{ji}v_i - \frac{kT_{,i}}{1-i\omega\tau_0}$ (freq. domain) & PARTIAL \\
GL & $J_i = \sigma_{ji}v_i + q_i$ (q from GL constitutive) & UNCERTAIN \\
GN-II & Entropy flux $p_i = k_{ij}\alpha_{,j}$ established; energy flux open & UNCERTAIN \\
GN-III & Mixed type I/II flux; open & UNCERTAIN \\
DPL & $J_i = \sigma_{ji}v_i - kT_{,i}\frac{1-i\omega\tau_T}{1-i\omega\tau_q}$ (freq. domain) & PARTIAL \\
TPL & Frequency-domain form formally derived; time domain open & UNCERTAIN \\
MGT & Frequency-domain reduces to classical; time domain open & UNCERTAIN \\
MDD & Convolution-form flux; kernel-dependent & UNCERTAIN \\
Nonlocal & Integral-form flux; divergence theorem modification needed & UNCERTAIN \\
Piezoelectric & $J_i^{\text{elec}} = -\phi\dot{D}_i$ & PARTIAL \\
\bottomrule
\end{tabularx}
\end{table}

\section{Recommendations}

\begin{enumerate}
\item Present only the classical flux $J_i = \sigma_{ji}v_i + q_i$ as canonical (verified).
\item Present frequency-domain forms for LS and DPL as theory-specific expressions with clearly stated assumptions.
\item Present all other formulations as open derivations with explicit assumptions.
\item Do not include any uncertain flux form in comparative tables as a canonical result.
\item Mark each uncertain flux as an open research question (OQ10.1--OQ10.8).
\end{enumerate}
